\chapter{Beispielimplementierung: Stromsensor SCT013}

	Die Implementierung der Klasse \PYTHON{CCurrentSensor} folgt der Struktur des Interfaces \PYTHON{ISensor} 
	und definiert die drei Pflichtmethoden:
	
	\begin{itemize}
		\item \PYTHON{init()}:
		
		Initialisiert den Sensor, setzt die Auflösung des ADC und konfiguriert die interne EmonLib-Instanz. 
		Nach erfolgreicher Initialisierung wird das Flag \PYTHON{initialized\_} aktiviert. 
		Diese Methode führt keine Messung durch, sondern bereitet den Sensor auf den Betrieb vor. 
		
		\textbf{Rückgabewert:} 
		
		\begin{itemize}
			\item \PYTHON{SensorError::SUCCESS} – die Initialisierung wurde erfolgreich abgeschlossen. 
		\end{itemize}
		
		\item \PYTHON{update()}:
		Führt eine neue RMS-Strommessung durch. Dazu werden mehrere ADC-Werte eingelesen, mit \PYTHON{EmonLib::calcIrms()} ausgewertet und das Ergebnis validiert. 
		Der gemessene Effektivwert wird anschlie{\ss}end in \PYTHON{lastReadVal\_} gespeichert.
		
		\textbf{Rückgabewerte:}
		
		\begin{itemize}
			\item \PYTHON{SensorError::SUCCESS} – die Messung wurde erfolgreich durchgeführt und der Wert gespeichert. 
			\item \PYTHON{SensorError::NOT\_INITIALIZED} – die Methode wurde aufgerufen, bevor \PYTHON{init()} ausgeführt wurde. 
			\item \PYTHON{SensorError::READ\_ERROR} – das Messergebnis war ungültig (z.\,B. \PYTHON{NaN} oder ein negativer Wert), 
			daher wurde der Wert verworfen. 
		\end{itemize}
		
		\item \PYTHON{getValue()}:
		
		Gibt den zuletzt gespeicherten Stromwert über eine Referenzvariable aus, 
		ohne eine neue Messung auszuführen. 
		Diese Methode dient ausschlie{\ss}lich dem Abruf des letzten gültigen Messwertes. 
		
		\textbf{Rückgabewerte:}

		\begin{itemize}
			\item \PYTHON{SensorError::SUCCESS} – ein gültiger Messwert liegt vor und wurde erfolgreich zurückgegeben. 
			\item \PYTHON{SensorError::NOT\_INITIALIZED} – der Sensor wurde noch nicht initialisiert, 
			daher ist kein gültiger Messwert verfügbar. 
		\end{itemize}
	\end{itemize}

{
  \captionof{code}{Klasse \PYTHON{CCurrentSensor}: Implementierung des Interfaces \PYTHON{ISensor} für RMS-Strommessung}
  \ArduinoExternal{}{../../Code/MKRWAN1310/SoftwareArchitecture/CCurrentSensor.h}
}
	
{
  \captionof{code}{Implementierung der Methoden des Stromsensors \PYTHON{CCurrentSensor}}
  \ArduinoExternal{}{../../Code/MKRWAN1310/SoftwareArchitecture/CCurrentSensor.cpp}
}
